\section{Methodology}
\label{methodology}

This section defines the proposed environment-aware elicitation workflow.
The method is intentionally lightweight: it uses role-based agent collaboration, environmental knowledge, and targeted human clarification to improve requirement context completeness.

\subsection{Framework Overview}
\label{method:overview}

The framework takes initial natural language requirements as input and produces context-enriched requirements plus a record of clarified environmental assumptions.
Its scope is limited to eliciting and completing environmental context.

\subsection{Multi-Agent Environment-Aware Workflow}
\label{method:workflow}

The workflow coordinates three analysis roles.

\subsubsection{Interviewer Agent}
\label{method:interviewer_agent}

The interviewer agent guides the elicitation dialogue, detects vague stakeholder statements, and decides when a clarification question should be raised.

\subsubsection{End-User Agent}
\label{method:end_user_agent}

The end-user agent examines requirements from the perspective of concrete user situations, including user goals, locations, interaction conditions, and task constraints.

\subsubsection{Environment Agent}
\label{method:environment_agent}

This agent identifies missing assumptions about deployment settings, devices, runtime conditions, connectivity, and operational constraints.

\subsection{Environment Information Detection}
\label{method:environment_detection}

This step checks whether each requirement contains enough environmental context for later analysis and implementation discussion.

\subsubsection{Deployment Environment}
\label{method:deployment_environment}

The workflow identifies assumptions about physical, organizational, network, and platform deployment settings.

\subsubsection{Device Constraints}
\label{method:device_constraints}

The workflow checks whether device type, interface capability, sensing availability, battery, performance, and accessibility constraints are specified when relevant.

\subsubsection{Runtime Conditions}
\label{method:runtime_conditions}

The workflow detects missing assumptions about connectivity, latency, workload, environmental changes, availability, and failure conditions.

\subsubsection{User Context Information}
\label{method:user_context}

The workflow captures user roles, usage scenarios, location, activity state, attention level, and interaction constraints.

\subsection{Knowledge-Assisted Context Understanding}
\label{method:knowledge_assisted}

Environmental knowledge is used to guide what context should be checked and what clarification questions should be asked.

\subsubsection{Environment Knowledge}
\label{method:environment_knowledge}

The method uses domain-relevant environment knowledge to identify likely deployment and runtime assumptions.

\subsubsection{Context Templates}
\label{method:context_templates}

Context templates define recurring fields such as deployment setting, device, user situation, runtime condition, and operational constraint.

\subsubsection{Scenario Patterns}
\label{method:scenario_patterns}

Scenario patterns help map vague requirements to concrete usage situations without introducing unrelated functionality.

\subsubsection{Prompt Strategies}
\label{method:prompt_strategies}

Prompt strategies constrain agents to ask environment-focused questions and to report uncertainty rather than invent missing context.

\subsection{Dynamic Environment Clarification}
\label{method:clarification}

Clarification is triggered when the workflow detects environmental gaps that materially affect requirement interpretation.

\subsubsection{Missing Environment Information}
\label{method:missing_environment}

The method records which environmental fields are absent, ambiguous, or inconsistent with the stated user scenario.

\subsubsection{Context Completion}
\label{method:context_completion}

The workflow updates requirements with stakeholder-confirmed context and separates confirmed assumptions from inferred assumptions.

\subsubsection{Scenario Refinement}
\label{method:scenario_refinement}

The workflow refines scenarios by linking requirements to concrete deployment, device, runtime, and user-context conditions.

\subsection{Human-in-the-Loop Interaction}
\label{method:hitl}

Human stakeholders review clarification questions and confirm or correct proposed environmental assumptions.
The interaction is designed to improve context completeness while keeping the clarification burden manageable.
